\documentclass[11pt,a4paper]{article}
\usepackage[T1]{fontenc}
\usepackage[utf8]{inputenc}
\usepackage{lmodern,microtype}
\usepackage[margin=27mm,headheight=14pt]{geometry}
\usepackage{amsmath,amssymb,amsthm,mathtools}
\usepackage{booktabs,array,longtable,enumitem}
\usepackage{xcolor,hyperref,fancyhdr,listings}
\definecolor{ink}{HTML}{173B58}
\definecolor{link}{HTML}{176B86}
\hypersetup{colorlinks=true,linkcolor=ink,urlcolor=link,citecolor=link,
 pdftitle={iPRA with inferred witnesses},pdfauthor={Research note prepared for Vladimir Reshetnikov}}
\pagestyle{fancy}\fancyhf{}
\fancyhead[L]{\small iPRA with inferred witnesses}
\fancyhead[R]{\small 11 November 2026}
\fancyfoot[C]{\thepage}
\setlength{\parskip}{4pt}
\setlength{\emergencystretch}{2em}
\newtheorem{theorem}{Theorem}[section]
\newtheorem{lemma}[theorem]{Lemma}
\newtheorem{proposition}[theorem]{Proposition}
\newtheorem{corollary}[theorem]{Corollary}
\theoremstyle{definition}\newtheorem{definition}[theorem]{Definition}
\theoremstyle{remark}\newtheorem{remark}[theorem]{Remark}
\newcommand{\ipra}{\mathsf{iPRA}}
\newcommand{\nf}{\operatorname{nf}}
\newcommand{\FV}{\operatorname{FV}}
\lstset{language=Python,basicstyle=\ttfamily\scriptsize,
 columns=fullflexible,keepspaces=true,breaklines=true,
 showstringspaces=false,frame=single,rulecolor=\color{black!15},
 keywordstyle=\color{ink},commentstyle=\color{black!50},
 aboveskip=10pt,belowskip=10pt}
\newcommand{\NewBytes}{1,990}

\title{\color{ink}iPRA with inferred witnesses\\[6pt]
 \large Three inference families, with equality supplied by a fixed numerical probe}
\author{Research note prepared for Vladimir Reshetnikov}
\date{11 November 2026}
\begin{document}
\maketitle
\begin{abstract}
We give a further reduction of the current ConciseIPRA presentation. Its
new elimination rule derives a numerical name from the zero and unit
instances of an existential matrix and specializes the opened clause at
that name. The certificate supplies neither a replacement term nor a proof
of the naming equality. A two-row numerical package recovers arbitrary
equality transport, while a separate substitution argument recovers ordinary
existential elimination even when substitution creates the naming pattern.
There is one fixed reflexivity clause and three inference families: clause
use, paired introduction, and opening with an inferred witness. Terms are
fully normalized scalar iterators; formulas are equations and a single
paired binder. We prove equality of the entire canonical clause consequence
relation with the established presentation. Composing with its arithmetic
bridge gives bi-interpretability with full first-order intuitionistic
primitive recursive arithmetic, with whole numerical domains, actual
equality, and identity numerical composites. A standalone Python checker,
a fuel-bounded variant, concrete certificates, and reproducible measurements
accompany the proof. The selected checker has \NewBytes{} normalized bytes.
The size improvement and the mathematical inference restriction are reported
separately. No absolute minimality theorem or new proof-assistant verification
is claimed.
\end{abstract}
\newpage
\tableofcontents
\newpage

\section{Result and the comparison that it supports}
\label{sec:result}

The construction retains the numerical sort, actual equality, and ordinary
reading of the paired binder from ConciseIPRA. Its change is local but
mathematical: an equality proof that previously had to be supplied at
witness discharge is replaced by information already contained in the
proved major formula. This is an admissible restriction of the equality
interface, with a proof of completeness for arbitrary finite contexts.
It does not rely on agreement on standard natural numbers alone.

\paragraph{The fixed baseline.}
All repository comparisons refer to the following commit of
\texttt{VladimirReshetnikov/ProveIt}:
\begin{center}\small\texttt{17898ea7e6a1300dfdcad71040d5daeed67916d4}\end{center}
The relevant directory is
\path{Logic/PrimitiveRecursiveArithmetic/ConciseIPRA}~\cite{repo}.
The README at that commit describes 67 filed packages. We also inspected
the six ZIP packages in its \texttt{docs/incoming} directory. In particular,
the incoming \texttt{witness\_discharge} proposal already has three inference
families and an explicit reflexivity clause~\cite{witness}; claiming that
rule count itself as new would be incorrect. Its elimination requires a
major proof, an equality proof, a branch proof, and an explicit replacement
term. The new rule needs only a major proof and a branch proof.

\paragraph{What is new.}
The fixed zero/unit test for an inferred witness, its use to specialize
the clause before opening, and its complete substitution and equality
recovery proofs are the new mathematical content. In particular, a syntactic
test need not be reflected by arbitrary substitution: preservation of its
positive case, together with strong induction on branch proofs, suffices.
Full normalization, optional-antecedent clauses, the reflexivity clause,
the single witness-selected introduction, numerical selector packages,
and the scalar arithmetic bridge are inherited and attributed.

\paragraph{What ``simpler'' means here.}
The proposal removes a proof premise and a supplied numerical argument
from the strongest inspected equality-discharge interface, and the delivered
complete implementation is smaller than the inspected explicitly typed
incoming checkers. The proof representation also changes from checking many
repeated conclusions to inferring each conclusion once. A matched control
quantifies the cost of the new naming test separately from that representation
change (Section~\ref{sec:implementation}). The repository's much shorter
compact codecs remain shorter. We do not call this an absolute source-size
minimum, and do not infer a lower bound from the absence of a smaller draft.

\paragraph{Strength and verification boundary.}
The intended iPRA is the full intuitionistic first-order theory with a
symbol for each finite primitive-recursive description, its defining
equations, the usual successor axioms, and induction for quantifier-free
formulas with parameters and retained assumptions. Induction is not added
for arbitrary quantified formulas. The new equivalence is a written
universal proof. The accompanying executions check the implementation and
specific certificates. Existing Lean sources certify the inherited bridge;
this package does not claim a new Lean proof of inferred opening.

% New mathematical fragment for synthesized naming AT OPENING.
% The enclosing report supplies standard theorem/lemma/corollary environments.
\section{Synthesizing a witness name when its clause is opened}
\label{sec:named-opening}

The new rule substitutes its inferred numerical name into
the opened clause, so its branch proves the desired outer conclusion
directly. The rule requires neither an equality proof, a supplied
replacement term, nor a separate branch-result formula. The main
completeness issue is admissibility of substitution. We resolve it before
recovering unrestricted existential-style elimination, using only two
guard-free forms of elimination that are immediately available.

\subsection{Complete presentation}

Numerical terms are full computation normal forms of
\[
 t::=x\mid0\mid S(t)\mid I(n,a;u.h)
\]
under
\[
 I(0,a;u.h)\longrightarrow a,\qquad
 I(Sn,a;u.h)\longrightarrow h[I(n,a;u.h)/u].
\]
Only $u$ in $h$ is bound. Write $N$ for normalization. Substitution is
simultaneous, capture avoiding, and followed by $N$. Inserting a term
does not recursively substitute into that inserted copy. The inherited
termination, confluence, and substitution-coherence results for $N$ are
part of the metatheoretic foundation and of the trusted checker boundary.

Formulas and optional-premise clauses are
\[
 A::=(s=t)\mid D_{x,y}B,\qquad
 [\Delta;y.B(y)],\quad |\Delta|\in\{0,1\},
 \quad y\notin\operatorname{FV}(\Delta).
\]
Each context is a finite list of clause occurrences containing the fixed
axiom occurrence $\mathcal R=[\,;y.y=y]$. An ordinary assumption $A$
means the unconditional constant clause $[\,;v.A]$. Put $H=(0=0)$,
$1=S0$, and $2=S(S0)$.

The intended ordinary readings are actual numerical equality,
\[
 (D_{x,y}B)^\circ=\exists x\bigl(B^\circ(x,0)\to
                                      \forall y\,B^\circ(x,y)\bigr),
\]
and $[\,;y.B]^\circ=\forall y B^\circ(y)$,
$[[A];y.B]^\circ=A^\circ\to\forall y B^\circ(y)$.
The three primitive inference families are the following.

\paragraph{Use.}
For a retained clause $[\Delta;y.B]\in\Gamma$,
\[
 \frac{\{\,\Gamma\vdash A:A\text{ is the antecedent occurrence in }
                                  \Delta\,\}}
      {\Gamma\vdash N(B(t))}.
 \tag{U}
\]
An empty antecedent requires no proof; a singleton requires one.
In particular, a use of $\mathcal R$ proves $t=t$.

\paragraph{Introduce.}
\[
 \frac{\Gamma,B(t,\eta)\vdash B(t,Sz)}
      {\Gamma\vdash D_{x,y}B(x,y)},\qquad
 \eta=\begin{cases}
 z,&B\text{ has equation root and }t\equiv0,\\
 0,&\text{otherwise}.
 \end{cases}
 \tag{I}
\]
The witness $t$ is canonical and $z$ is fresh for the retained context,
the witness, and all passive matrix parameters. Literal syntax, not
provable equality, determines the mode.

\paragraph{Open with a synthesized name.}
For a major $D_{x,y}B(x,y)$ choose a fresh parameter $b$ and put
\[
 B_b(y)=N(B(b,y)),\qquad
 \tau_B(b)=\begin{cases}
 r,&B_b(0)\equiv H\text{ and }B_b(1)\equiv(b=r),\\
 b,&\text{otherwise}.
 \end{cases}
 \tag{P}
\]
Define the substituted clause
\[
 \mathcal C^*_B(b)=
 [[N(B(\tau_B(b),0))];y.N(B(\tau_B(b),y))].
\]
Then the rule is
\[
 \frac{\Gamma\vdash D_{x,y}B(x,y)\qquad
       \Gamma,\mathcal C^*_B(b)\vdash A}
      {\Gamma\vdash A},
 \qquad b\notin\operatorname{FV}(\Gamma,A).
 \tag{O$^*$}
\]
The parameter $b$ is also fresh for the passive parameters of $B$.
The synthesized $r$ may contain $b$; the branch still lives in the scope
with $b$ available, and its conclusion must be in the outer scope. The
primitive rule has exactly two proof premises. There is no additional
computation, equality, hypothesis, induction, cut, or substitution rule.
Computation is part of canonical syntax, and reflexivity is the explicit
fixed axiom clause.

\begin{lemma}[Direct ordinary soundness]
\label{lem:opening-sound}
O$^*$ is sound in intuitionistic first-order equality logic over the scalar
computation equations.
\end{lemma}
\begin{proof}
Open the existential major at $b$, obtaining
$\mathcal C_B(b)=[[B_b(0)];y.B_b(y)]$. If P selects its first branch,
the zero instance is reflexive and clause use at $1$ gives $b=\tau_B(b)$.
Otherwise reflexivity gives the same equality because $\tau_B(b)=b$.
Ordinary equality transport converts the whole clause reading
$\mathcal C_B(b)^\circ$ to
$\mathcal C^*_B(b)^\circ$. Apply the supplied branch derivation to this
converted assumption. The conclusion contains no $b$, so ordinary
existential elimination closes. Equality transport and simultaneous
substitution are valid also for an equality $b=r(b)$.
\end{proof}

\subsection{The preliminary rules that do not need general substitution}

Weakening by clause occurrences, injective fresh renaming, and cut of an
ordinary formula assumption are admissible by direct induction on
derivations. For renaming, the injectivity and freshness preserve and
reflect the literal patterns in I and P; normalization is equivariant.
For ordinary-assumption cut, replace only uses of the selected
unconditional constant clause, retain all other clause occurrences, and
freshen local parameters before insertion. These arguments do not
substitute arbitrary numerical terms into a derivation.

\begin{lemma}[Two immediately ordinary forms of opening]
\label{lem:opening-safe}
O$^*$ is exactly ordinary paired elimination in each of the following cases:
\begin{enumerate}
 \item the witness coordinate $x$ is unused in the matrix $B(x,y)$;
 \item the matrix $B$ has compound, rather than equation, root.
\end{enumerate}
\end{lemma}
\begin{proof}
In the first case $B_b(1)$ contains no $b$, so its left-hand term cannot
literally be the fresh parameter $b$. In the second case $B_b(0)$ still
has compound root, because term normalization does not erase formula
binders, and therefore cannot be $H$. In either case P returns $b$.
\end{proof}

For unused $w$ write $q_yA(y)=D_{w,y}A(y)$.
\begin{lemma}[Closures and numerical cases]
\label{lem:opening-cases}
Before using general ordinary elimination or equality transport, the
following rules are available:
\[
 \frac{\Gamma,A(0)\vdash A(Sz)}{\Gamma\vdash q_yA(y)},
 \qquad
 \frac{\Gamma\vdash q_yA(y)\qquad\Gamma\vdash A(0)}
      {\Gamma\vdash A(t)}.
\]
Consequently $\Gamma\vdash A(0)$ and $\Gamma\vdash A(Sz)$, with $z$
fresh, imply $\Gamma\vdash A(t)$ for every well-scoped $t$.
\end{lemma}
\begin{proof}
Introduce the unused witness of $q$ at $1$, forcing I's case mode.
For specialization, use the unused-witness case of
Lemma~\ref{lem:opening-safe}, which opens the clause
$[[A(0)];y.A(y)]$; U at $t$ with the given base proof gives the desired
outer conclusion. The numerical-case rule follows by weakening the
successor proof with the unused assumption $A(0)$ and composing the
two constructions.
\end{proof}

\begin{lemma}[Unrestricted case introduction]
\label{lem:opening-case-repair}
For every canonical matrix $B$ and every canonical witness $t$, the rule
\[
 \frac{\Gamma,B(t,0)\vdash B(t,Sz)}
      {\Gamma\vdash D_{x,y}B(x,y)}
\]
is derivable, with the same freshness condition on $z$ as in I. Its
derivation uses only U, I, and the two ordinary opening forms of
Lemma~\ref{lem:opening-safe}.
\end{lemma}
\begin{proof}
The primitive I already suffices unless $B$ has equation root and $t=0$.
Write the remaining matrix as $E(x,y)$. Put
\[
 C(v,a,c)=N(I(v,a;u.c)),\quad u\notin\operatorname{FV}(c),
\]
and for an equation $E=(s=t)$ define
\[
 G_vE=(C(v,0,s)=C(v,0,t)),\qquad L(E)=q_vG_vE.
\]
The literal endpoints are $G_0E=H$ and $G_{Sw}E=E$.
Lemma~\ref{lem:opening-cases} gives $E\vdash L(E)$ by closure
introduction and $L(E)\vdash E$ by specialization at $1$ and reflexivity
of $H$. Only unused-witness opening is involved.

Let $A(x,y)=L(E(x,y))$. Under $A(0,0)$, unwrap to $E(0,0)$, insert the
supplied premise by ordinary-assumption cut, and wrap $E(0,Sz)$ as
$A(0,Sz)$. Since $A$ is compound, primitive I at zero uses cases and
gives $D_{x,y}A(x,y)$. Open this major at fresh $c$ using the
compound-matrix case of Lemma~\ref{lem:opening-safe}. To introduce
$D_{x,y}E(x,y)$ at witness $c$, assume $E(c,0)$, wrap it to $A(c,0)$,
use the opened clause at $Sz$, and unwrap $A(c,Sz)$ to $E(c,Sz)$.
The fresh witness $c$ is not literal zero, so I uses cases. The resulting
$D_{x,y}E(x,y)$ is independent of $c$ and can be exported. No general
substitution theorem and no unrestricted opening rule was used.
\end{proof}

\subsection{The substitution theorem and its noncircular proof}

A simultaneous numerical substitution may send a parameter to a term
containing that same parameter. Images are inserted once. Only the finitely
many parameters appearing in a given derivation matter. As usual, bound
variables and eigenparameters are freshened away from the images before
the corresponding substitution is carried out.

\begin{theorem}[Admissibility of arbitrary numerical substitution]
\label{thm:opening-substitution}
If $\Gamma\vdash A$ and $\sigma$ is any well-scoped simultaneous
numerical substitution, then
\[
 N(\Gamma\sigma)\vdash N(A\sigma).
\]
This is an effective transformation of finite derivations. It is not a
claim that substitution preserves every primitive rule instance.
\end{theorem}
\begin{proof}
Use strong induction on the original derivation, with the induction
assertion quantified over \emph{all} simultaneous substitutions. This
quantification is essential: in one opening case the branch receives a
different combined substitution from the major. The induction hypothesis
is applied to original proper subderivations, never to an already
translated derivation of potentially larger size.

For U, substitute its selected clause occurrence and its zero or one
premises, then use normalization--substitution coherence. The distinguished
axiom $\mathcal R$ is unchanged because it has no free numerical parameter.
For I, first freshen $z$ away from all images of $\sigma$. A progress
instance stays progressive: its equation root and literal-zero witness
are preserved. A case instance stays a case instance unless an equation
matrix has a previously nonzero witness whose substituted normal form
is zero. In that exceptional case, the induction hypothesis gives exactly
the premise of the derived unrestricted case rule in
Lemma~\ref{lem:opening-case-repair}; apply that rule. That lemma was proved
before the present substitution theorem.

It remains to treat O$^*$. Freshen its parameter $b$ away from the passive
data and from the finitely many substitution images. Write
$B'=N(B\sigma)$ and $\tau'=\tau_{B'}(b)$. Translate the major by the
induction hypothesis. There are two cases according to the \emph{old}
probe.

If the old probe selected $r$, its two literal identities were
$B_b(0)=H$ and $B_b(1)=(b=r)$. Applying $\sigma$ while leaving $b$ fixed
preserves these identities after normalization: $H$ is closed and $b$
is untouched. Formula roots are preserved. Therefore the new probe is
also successful and $\tau'=N(r\sigma)$ literally. Apply the induction
hypothesis to the original branch using $\sigma$ with $b$ fixed. Its
substituted clause is, by coherence,
\[
 [[B'(\tau',0)];y.B'(\tau',y)],
\]
which is precisely the clause needed by the new O$^*$ instance.

If the old probe used its fallback, its branch clause was the ordinary
$[[B(b,0)];y.B(b,y)]$. The new probe may remain false or may become true.
In \emph{both} cases apply the induction hypothesis once to the original
branch with the combined simultaneous substitution that agrees with
$\sigma$ on the outer parameters and sends $b$ to $\tau'$. Its image
$\tau'$ is inserted once, including when it contains $b$. The resulting
clause is exactly
$[[B'(\tau',0)];y.B'(\tau',y)]$. The outer context and branch conclusion
become $N(\Gamma\sigma)$ and $N(A\sigma)$ because neither originally
contained $b$, and the freshening ensured that no outer image introduces
$b$. This supplies the new branch whether or not the guard was newly
created. Close the new O$^*$ instance.

This exhausts the rules. The only proof transformation invoked outside
the induction hypotheses was the already established case repair.
\end{proof}

\begin{corollary}[Unrestricted ordinary paired elimination]
\label{cor:opening-ordinary}
For arbitrary canonical $B$, the rule
\[
 \frac{\Gamma\vdash D_{x,y}B(x,y)\qquad
       \Gamma,[[B(b,0)];y.B(b,y)]\vdash A}
      {\Gamma\vdash A},\qquad
 b\notin\operatorname{FV}(\Gamma,A),
\]
is admissible, where $b$ is also fresh for the passive parameters of $B$.
\end{corollary}
\begin{proof}
Apply Theorem~\ref{thm:opening-substitution} to the supplied branch,
with the simultaneous substitution $b\mapsto\tau_B(b)$ and all outer
parameters fixed. Since $\Gamma,A$ are $b$-free, the result is exactly
the branch of O$^*$ with the same outer conclusion $A$.
The major is unchanged. This proves general ordinary elimination only
\emph{after} the substitution theorem, avoiding circularity.
\end{proof}

\subsection{Recovering every equality rule}

With $C$ as above define
\[
 P(y)=N(I(y,0;u.C(u,1,0))).
\]
Computation gives $P(1)=1$ and $P(2)=0$ literally; no parity induction is
needed.

\begin{theorem}[Uncrossed transport]
\label{thm:opening-transport}
For every fixed numerical template $h(w)$,
\[
 \Gamma\vdash s=t,\quad\Gamma\vdash N(h(s))=r
 \quad\Longrightarrow\quad\Gamma\vdash N(h(t))=r.
\]
\end{theorem}
\begin{proof}
Use the canonical equation matrix
\[
 B(x,y)=
 \bigl(C(y,0,C(P(y),N(h(x)),x))
      =C(y,0,C(P(y),r,t))\bigr).
\]
Its zero, first, and second rows are respectively
\[
 H,\qquad x=t,\qquad N(h(x))=r.
\]
At witness $s$, numerical cases for
$C(v,N(h(s)),s)=C(v,r,t)$, with the two supplied premises as its zero
and successor cases, give the successor row at $v=P(Sz)$. This ignores
I's local assumption, so I introduces $D_{x,y}B(x,y)$ in either mode.
Open that major. Its probe returns $t$, so its branch clause is
\[
 [[H];y.B(t,y)].
\]
Reflexivity proves $H$, and U at $2$ proves $N(h(t))=r$ directly.
That equation is in the outer scope, so O$^*$ closes. The numerical-case
construction was derived before equality, and the transport construction
does not require a general substitution of the supplied equality proofs.
\end{proof}

Symmetry follows by taking $h(w)=w$ and the reflexive minor $s=s$;
crossed transport follows by symmetry of the output. Transitivity and
term congruence follow in the usual finite equational constructions.
Replacing both sides of an equation uses two one-term transports and
symmetry.

Full formula transport follows by external induction on formula depth.
For a compound template $D_{x,y}B(v,x,y)$, open the source formula by
Corollary~\ref{cor:opening-ordinary} at a fresh $c$, and introduce the
target formula at witness $c$, which forces case mode. Under its zero
assumption, transport backward in the smaller matrix, specialize the
source clause at $Sz$, and transport forward in the smaller matrix.
Close introduction and ordinary elimination. This uses no object-theory
induction on quantified formulas.

\subsection{Identity-language equivalence and both interpretation composites}

Let $\mathcal T$ be the inherited fully normalized clause calculus with
the same syntax, Ax, U and I, together with ordinary paired elimination
and crossed one-term transport. Let $\mathcal O$ be the new three-rule
calculus U, I, O$^*$.

\begin{theorem}[Exact equivalence on arbitrary clause sequents]
\label{thm:opening-equivalence}
For every canonical finite clause context $\Gamma$ containing Ax and
every canonical formula $A$,
\[
 \Gamma\vdash_{\mathcal O}A
 \quad\Longleftrightarrow\quad
 \Gamma\vdash_{\mathcal T}A.
\]
Both translations are effective and leave the full end sequent literally
unchanged.
\end{theorem}
\begin{proof}
From $\mathcal T$ to $\mathcal O$, translate proof trees recursively,
replacing ordinary elimination by Corollary~\ref{cor:opening-ordinary}
and crossed transport by Theorem~\ref{thm:opening-transport} and derived
symmetry. U and I are shared. The preliminary admissibility theorems were
proved internally in $\mathcal O$ and do not rely on this translation.

Conversely, $\mathcal T$ has full formula transport by the same
formula-depth argument. Translate an O$^*$ major and open it ordinarily
at $b$, retaining its clause $\mathcal C_B(b)$. Derive
$b=\tau_B(b)$ by reflexivity and U at $1$ when the probe succeeds, or
by reflexivity alone otherwise. In this opened context derive
$q_y B(\tau_B(b),y)$ as follows. Under its zero assumption
$B(\tau_B(b),0)$, transport backward along the derived equality to
$B(b,0)$; use the old clause at $Sz$; transport forward to
$B(\tau_B(b),Sz)$; and introduce the unused $q$ witness at $1$.
Ordinary elimination of that $q$ opens exactly
$\mathcal C^*_B(b)$. Insert the translated O$^*$ branch, weakened by
the still-retained old clause, and close this unused-witness elimination.
The branch conclusion is free of $b$, so the outer ordinary elimination
also closes. All operations are on finite proof trees, with fresh
parameters chosen before insertion.
\end{proof}

\begin{corollary}[Full iPRA and the identity numerical composites]
\label{cor:opening-ipra}
The new calculus is equiconsistent and bi-interpretable with full
first-order intuitionistic primitive recursive arithmetic. Both
interpretations use the whole numerical domains and actual equality;
both numerical composite isomorphisms have graph $x=y$.
\end{corollary}
\begin{proof}
Compose Theorem~\ref{thm:opening-equivalence} with the inherited canonical
clause/native-scalar comparison and the inherited scalar/iPRA
bi-interpretation. The new comparison is identity on every numerical term
and formula, so it changes neither numerical map nor the formula
translations. Each new composite therefore consists of identity legs
and an inherited composite already provably equivalent to the original
formula, with numerical graph $x=y$. The designated contradiction
$0=S0$ is unchanged, giving equiconsistency as a special case.
\end{proof}

The new universal metatheorems are
Theorems~\ref{thm:opening-substitution} and
\ref{thm:opening-equivalence}. Their dependency order is: structural
renaming and ordinary-assumption cut; the two guard-free openings;
closure rules; case repair; strong-induction substitution; unrestricted
ordinary opening; equality; exact comparison. The iPRA arithmetic
compiler and normalization foundation remain explicitly inherited
results. No newly proof-assistant-checked status or primitive-recursive
bound on full normalization is asserted.

\section{The normalization layer is part of the theory}
\label{sec:normalization}

Moving computation into syntax removes two displayed computation axioms,
but it does not remove their implementation. The checker includes full
normalization for exactly
\[
 I(0,a;u.h)\longrightarrow a,
 \qquad I(Sn,a;u.h)\longrightarrow h[I(n,a;u.h)/u],
\]
closed under every term context, including iterator bodies. There is no
reduction of formula constructors. The following argument records the
inherited foundation~\cite{normalization,witness} so that the new probe
does not depend on an unstated computation oracle.

\begin{lemma}[Closure under iteration]
Suppose $n$ and $a$ are strongly normalizing, and $h[v/u]$ is strongly
normalizing for every strongly normalizing $v$. Then $I(n,a;u.h)$ is
strongly normalizing.
\end{lemma}
\begin{proof}
For a strongly normalizing term $s$, let $\nu(s)$ be the maximum length
of a reduction from it. It exists because the reduction tree is finite
branching and has no infinite path. In particular $h$ is strongly
normalizing by taking $v=u$. Induct lexicographically on
\[
 (\nu(n),\ |n|,\ \nu(a)+\nu(h)).
\]
A reduction inside $n$ decreases the first component. A reduction inside
$a$ or $h$ decreases the last; the universal hypothesis on $h$ survives a
reduction in $h$, because capture-free substitution preserves reduction.
The zero root reduct is $a$. At a successor root $n=Sp$, the predecessor
$p$ has the same reduction height and smaller size. Thus $I(p,a;u.h)$
is strongly normalizing by induction, and its substitution into $h$ is
strongly normalizing by the hypothesis on $h$. Every immediate reduct
is strongly normalizing, which proves the claim.
\end{proof}

\begin{theorem}[Unique normal forms and substitution coherence]
Every finite raw term has a unique normal form up to bound-variable
renaming. With de Bruijn indices the normal form is unique literally.
For every finite simultaneous substitution $\sigma$,
\[
 N(t\sigma)=N\bigl(N(t)[x\mapsto N(\sigma x)]\bigr).
 \tag{Coherence}
\]
The same identity holds for formulas and clauses, and normalized
substitutions compose in the corresponding capture-avoiding way.
\end{theorem}
\begin{proof}
Structural induction on $t$ proves the stronger statement that every
substitution of strongly normalizing terms into $t$ yields a strongly
normalizing term. The iterator case uses the preceding lemma; the
induction hypothesis on its body applies with every additional strongly
normalizing assignment to the state variable. Identity substitution gives
strong normalization of all raw terms.

For local confluence, disjoint contractions commute, and the zero and
successor root patterns do not overlap. If a root contraction copies a
subterm containing the other contraction, perform that contraction in
each of its finitely many residual copies. If it erases the subterm,
no residual contraction is required. Capture-free substitution supplies
the common reduct in the iterator-body case. Induction on reduction
height now proves confluence. A substituted reduction of $t$ and the
reductions of the finitely many substituted images give convertible
terms on the two sides of Coherence; their unique normal forms agree.
Formula and clause constructors are preserved, so the same proof extends
componentwise to them.
\end{proof}

\paragraph{Connection with raw syntax.}
In the raw scalar calculus each reduction is a provable equation by a
computation axiom and congruence. Therefore every fixed term is provably
equal to its normal form, and every formula is provably equivalent to
its normalized formula. In the opposite direction the normal forms of
both computation axioms are reflexive equations. Transport and clause
use commute with normalization by Coherence. The only subtle introduction
case is a nonzero-looking raw witness that becomes literal zero; the
compound-envelope case repair supplied above handles that instance.
These observations give the canonical/raw comparison used in
Corollary~\ref{cor:opening-ipra}. The raw formulas and their normal forms
represent the same ordinary first-order statements up to provable
equivalence.

\paragraph{No uniform numerical evaluator is postulated.}
For each fixed finite iterator description, iPRA can interpret its
operation by primitive recursion. That does not give a primitive-recursive
function uniformly evaluating all primitive-recursive descriptions. In
fact such a uniform evaluator would contradict the familiar diagonal
construction $n\mapsto U(n,n)+1$. The normalization theorem above is a
metatheorem about all finite terms; it is not an extra totality axiom
inside iPRA. The complete checker includes the normalization code, and
the fuelled checker supplies a bounded certificate relation without
claiming a uniform primitive-recursive normalization bound.

\paragraph{Why canonical syntax is checked, not merely evaluated.}
Input terms are required to be normal forms. The syntax traversal checks
their constructor shapes and variable scopes, computes their normal
forms, and requires the result to be identical to the input. Rule-created
substitution instances may expose new redexes; those are normalized
before use. In particular the probe uses normalized zero and unit
instances, not a textual search through an unnormalized matrix. It can
therefore discover a name that becomes visible only after computation.

\section{The full intuitionistic theory and both composites}
\label{sec:bridge}

Exact equivalence of clause sequents is stronger than equality of closed
equational theorems. Here we spell out how it transfers the complete
iPRA interpretation, including its logical operations and both round
trips. This section summarizes inherited results~\cite{repo,lean,arithmetic};
their detailed written arithmetic proofs are included in \texttt{reference}.

\subsection{Ordinary realization}
Let $T_I$ be intuitionistic first-order logic with actual equality,
genuine falsity $\bot$, and the staged scalar-iterator language. Its
nonlogical axioms and schemas are the two iterator computations,
\[
 \forall x\,(x=0\vee\exists y\,x=Sy),
 \qquad (0=S0)\to\bot,
\]
and induction for every literal equation $E(y,\vec p)$:
\[
 \bigl(E(0,\vec p)\wedge
       \forall y(E(y,\vec p)\to E(Sy,\vec p))\bigr)
       \to\forall y E(y,\vec p).
\]
Equivalently this is the usual contextual induction rule, with the
induction variable fresh for retained assumptions. A staged iterator
symbol $i_h(n,a,\vec p)$ is introduced for each fixed earlier-stage
description $h(u,\vec p)$; no object-language variable ranges over
arbitrary function descriptions. The numerical domain is the whole
domain and interpreted equality is equality itself.

Use has its ordinary conditional-universal reading. Introduce is sound
by numerical coverage in case mode and equation induction in progress
mode. Opening is sound by Lemma~\ref{lem:opening-sound}. The explicit
falsity bridge must not be omitted from $T_I$: positive arithmetic
equations alone admit a one-element structure. In the native calculus
the designated contradiction is the equation $0=S0$.

\subsection{All logical operations as macros}
Write $C(v,a,b)=N(I(v,a;u.b))$ with $u$ unused, and let
$q_y A(y)=D_{w,y}A(y)$ for unused $w$. A gate on formulas is defined by
external recursion on their finite syntax:
\[
 \begin{aligned}
 G_v(s=t)&=(C(v,0,s)=C(v,0,t)),\\
 G_v(D_{x,y}B)&=D_{x,y}G_vB.
 \end{aligned}
\]
All bound variables are renamed away from $v$. Its endpoints satisfy
$\vdash G_0A$ and $G_{Sz}A\dashv\vdash A$; also
$A\dashv\vdash q_yG_yA$ when $y$ is fresh. These are derivability
statements, including for nested binders.

For equations define the selector
\[
 M_v(s=t,u=w)=(C(v,s,u)=C(v,t,w)).
\]
For compound formulas align fresh paired binders and define
\[
 \begin{aligned}
 M_v(D_{x,y}A,D_{x,y}B)&=D_{x,y}M_v(A,B),\\
 M_v(E,D_{x,y}B)&=D_{x,y}M_v(G_yE,B),\\
 M_v(D_{x,y}A,E)&=D_{x,y}M_v(A,G_yE),
 \end{aligned}
\]
where $E$ is an equation. Each recursive call lowers the maximum binder
depth. Gate padding and binder congruence give
$M_0(A,B)\dashv\vdash A$ and $M_{Sz}(A,B)\dashv\vdash B$.
With $P(y)=N(I(y,0;u.C(u,1,0)))$, the following complete inventory is
available:
\[
\begin{aligned}
 \bot_*&=(0=S0),\\
 A\Rightarrow_*B&=q_yM_y(A,B),\\
 A\wedge_*B&=q_yG_yM_{P(y)}(A,B),\\
 \exists_*x\,A(x)&=D_{x,y}G_yA(x),\\
 A\vee_*B&=\exists_*v\,M_v(A,B),\\
 \forall_*x\,A(x)&=A(0)\wedge_*q_yG_yA(y).
\end{aligned}
\tag{Logical macros}
\]
Freshness is understood in every line. The last line intentionally
supplies the zero instance separately: the gated closure supplies all
successor instances.

Implication introduction and elimination use selector endpoints and the
closure rules. Existential introduction and elimination use gate
endpoints. Conjunction is projected at $1$ and $2$, using $P(1)=1$,
$P(2)=0$; introduction uses numerical cases on its selector. Disjunction
uses control witnesses $0$ and $1$, and its elimination uses numerical
cases. For universals, numerical coverage combines the zero conjunct
with the successor instances in the closure. Under $0=S0$, equality
transport identifies $C(0,s,t)$ and $C(S0,s,t)$, proving every numerical
equation; external induction on formula depth and case introduction
then prove every formula. Thus all usual intuitionistic first-order
inference rules, including explosion, are recovered.

These macros do not add unrestricted induction for quantified formulas.
Primitive progress remains restricted to equations. Equation induction itself is obtained
by introducing a closure at an unused zero witness, opening it, and
supplying its base. The explicit repaired case rule handles the case
introduction instances needed above.

\subsection{The numerical arithmetic bridge}
The inherited scalar proof has a necessary dependency order. First it
constructs a scalar predecessor and proves $\operatorname{pred}(0)=0$
and $\operatorname{pred}(Sn)=n$ using equation induction and coverage.
It then constructs truncated subtraction and a characteristic term
for actual equality on the entire domain. Compiling a fixed
quantifier-free formula into its characteristic equation gives
quantifier-free induction. Only after this stage does it construct
pairing and compile arbitrary primitive recursion. These steps are
proved internally, rather than checked only on standard numerals.

For concreteness, once predecessor, subtraction, and quantifier-free
induction are available, a bounded square-root construction proves
pairing and projections $\pi,L,H$ satisfying
\[
 L(\pi(a,b))=a,\qquad H(\pi(a,b))=b.
\]
One choice is $\pi(a,b)=(a+b)^2+a$; bounded search for the integer
square root supplies its projections. Surjectivity of $\pi$ is
unnecessary: it stores a recursion state and does not replace the
interpretation domain by a coding range.

Given a fixed primitive recursion
\[
 f(0,\vec p)=g(\vec p),\qquad
 f(Sn,\vec p)=h(n,f(n,\vec p),\vec p),
\]
and already translated $g^*,h^*$, set
\[
 \begin{aligned}
 W(u,\vec p)&=\pi(SL(u),h^*(L(u),H(u),\vec p)),\\
 V(n,\vec p)&=I(n,\pi(0,g^*(\vec p));u.W(u,\vec p)),\\
 f^*(n,\vec p)&=H(V(n,\vec p)).
 \end{aligned}
\]
Equation induction proves $L(V(n,\vec p))=n$ and the defining equations
for $f^*$. Initial functions and composition translate directly.
In the reverse direction, every fixed scalar iterator is an ordinary
primitive recursion whose step ignores the recursion counter.

\subsection{Why this is bi-interpretability}
For the numerical composites, internal equation induction proves
uniqueness of each fixed iterator or primitive recursion satisfying
its defining equations. External induction over a finite function
description therefore identifies its double translation with the
original operation at every internal argument. The identity graph
$x=y$ is consequently a definable total bijection preserving every
operation in both composite interpretations. It is not a retract,
quotient, or restriction to standard numerals.

For the logical composites, the essential native identity is
\[
 D_{x,y}B(x,y)\dashv\vdash
 \exists_*x\bigl(B(x,0)\Rightarrow_*\forall_*y B(x,y)\bigr).
\]
Forward, ordinary opening supplies the conditional-universal clause;
introduce the three macros and close the opening. Backward, eliminate
the existential macro, introduce $D$ at its witness using the repaired
case rule, and obtain the required successor instance from implication
and universal elimination. The endpoint proofs above give the round
trips for all ordinary connectives and quantifiers. External induction
on formulas finishes both composites. Optional clauses correspond to
their unconditional or conditional universal readings.

These are exactly the inherited bridges recorded by
\texttt{biInterp} and the four faithfulness and formula round-trip theorems in
\texttt{NativeBiInterpretation.lean}~\cite{lean}. The new theorem
preserves every end sequent of the canonical clause presentation,
so the interpretations and the identity numerical graphs transfer
unchanged. In particular the new calculus derives $0=S0$ if and
only if iPRA derives $\bot$. This establishes equiconsistency; it
does not assert that either theory proves its own consistency.

\section{Complete implementation and measured simplification}
\label{sec:implementation}

\subsection{Input contract and trusted code}
The primary program reads a JSON object with exactly two fields,
\texttt{goal} and \texttt{proof}. The goal is a closed canonical formula.
The checker prints \texttt{ACCEPT} only after inferring the proof's
conclusion, validating its syntax and scope, and comparing it with the
independently validated goal. An exception or nonzero process exit rejects.
No Python \texttt{assert} is used in the kernel, so optimized execution
does not disable a logical check.

Zero is integer \texttt{0}; negative integers are free parameter levels;
positive integers are bound indices, with index one innermost. Terms
\texttt{[0,t]} and \texttt{[n,a,h]} encode successor and iteration, with
one bound index added only inside \texttt{h}. A two-element formula
\texttt{[s,t]} is an equation; \texttt{[B]} binds two indices in its
matrix, with the counter at index one and the witness at index two.
The sorted syntax traversal distinguishes a successor from an equation.
Booleans, floats, out-of-scope indices, and noncanonical terms are rejected.

The instruction formats are
\[
\begin{array}{ll}
 \texttt{["u",j,t,p]}&\text{use clause }j\text{ at }t,\\
 \texttt{["i",B,t,p]}&\text{introduce }[B]\text{ at witness }t,\\
 \texttt{["e",p,q]}&\text{infer a major from }p\text{ and check branch }q.
\end{array}
\]
For unconditional clause use, the fourth payload is unused; the examples
use JSON \texttt{null}. This is a declared certificate convention, not
an unvalidated proof premise. Conditional use must infer exactly the
retained antecedent. The context initially contains only
\texttt{([], [1,1])}, the reflexivity clause. Every I or O$^*$ branch adds
one clause and one parameter, so the context length determines the next
fresh level without a separate numerical-scope argument.

Every successful recursive call validates its inferred formula in that
call's outer scope. Thus opening cannot export its fresh parameter,
including when the synthesized term contains it. Introduction returns
the unchanged annotated matrix and validates that matrix at the end;
malformed syntax temporarily erased in a computed branch cannot survive
this final check. The independent validation includes such erased
malformation examples.

\paragraph{The retained-antecedent invariant.}
The implementation computes the original antecedent $u=N(B(b,0))$ and
substitutes the synthesized name into the body. It need not recompute
the antecedent. If the probe succeeds, $u$ is literally $H$, and
substitution leaves it unchanged. If it fails, the substitution is the
identity. Therefore the stored pair is exactly the clause stipulated
by O$^*$ in both cases. This identity accounts for part of the matched
opening-control comparison below.

\subsection{Source measurements}
The measurement convention follows the repository: remove a module's
documentation string, parse the remaining source, normalize it with
CPython 3.12.14 \texttt{ast.unparse}, and append one newline. Count UTF-8
bytes, significant Python tokens (excluding comments, indentation,
newlines and end markers), and all nodes from \texttt{ast.walk}.
The gzip column uses level 9, a zero timestamp, and these normalized
bytes. Gzip is an informational metric, not an optimization objective.
All trusted syntax, normalization, binding, scope, proof, and JSON
handling code is included. No untrusted proof-construction code is
required for acceptance.

\begin{table}[htbp]
\centering\small
\begin{tabular}{@{}lrrrr@{}}
\toprule
Complete program & Bytes & Tokens & AST & Gzip 9\\
\midrule
Inferred opening (selected) & 1,990 & 683 & 771 & 819\\
Matched equality opening & 1,998 & 676 & 765 & 812\\
Matched equality output & 1,983 & 663 & 751 & 812\\
Incoming witness discharge & 2,273 & 732 & 813 & 903\\
Incoming uniform binding & 2,238 & 693 & 764 & 866\\
Incoming clause contexts & 2,282 & 666 & 770 & 774\\
Inferred opening with fuel & 2,203 & 728 & 826 & 885\\
Incoming witness with fuel & 2,486 & 777 & 868 & 967\\
\bottomrule
\end{tabular}
\caption{Complete programs under one reproducible source-measurement convention.}
\label{tab:metrics}
\end{table}


Against the incoming explicit-equality discharge implementation, the
selected program saves 283 normalized bytes (12.45\%), 49 significant
tokens (6.69\%), and 42 AST nodes (5.17\%); gzip falls by 84 bytes.
This comparison includes both the new primitive interface and the
change to a forward checker that infers conclusions once. It must not
be attributed entirely to the mathematical rule restriction.

The matched explicit-equality \emph{opening} control uses the same
syntax traversal, normalizer, context convention, and forward interface.
Its opening accepts an equality proof $b=t$, specializes the clause at
$t$, and checks the branch. The new checker saves eight bytes, while
using seven more significant tokens, six more AST nodes, and seven more
gzip bytes. The matched explicit-equality \emph{output} control instead
keeps the unspecialized clause and substitutes into the branch's result;
it is seven normalized bytes shorter than the selected checker.
These mixed results are recorded rather than concealed: removing a
proof premise does not automatically shorten all implementations.

The selected checker is shorter in normalized bytes than the six
inspected explicitly typed incoming checkers, but not in every coordinate.
For example, \texttt{uniform\_binding} has seven fewer AST nodes, and
\texttt{clause\_contexts} has fewer tokens and one fewer AST node.
The repository's 516-byte compact codec and 788-byte validated-public-goal
implementation remain substantially shorter~\cite{repo}. Their
representation and internal validation contracts differ, so they are
not matched controls; their smaller sizes nonetheless preclude a
global source-size record claim for this package.

\subsection{The fuel-bounded certificate relation}
The second standalone checker adds an integer \texttt{fuel} field. It
decrements fuel on every syntax/normalization traversal entry and proof
call, rejecting before the counter becomes negative. Its full size is
included in the table. For each accepted unbounded certificate, its
finite computation uses some finite amount of fuel; conversely every
fuelled acceptance follows the same checking computation. The two
idealized checkers therefore accept proofs of exactly the same formulas
when sufficient resources are supplied.

With input size and fuel fixed, the number of recursive entries is
bounded. Each entry performs finite operations on syntax whose growth
has a primitive-recursive bound in those two inputs. The fuelled
certificate relation is consequently primitive recursive. This does
not provide a primitive-recursive bound on the fuel needed to normalize
all descriptions, and does not claim that iPRA proves the uniform
normalizer total. As with the baseline Python programs, implementation
stack, memory, and parser limits can reject a large mathematically
valid certificate; such rejection never licenses acceptance of an
invalid one.

\subsection{Executed evidence}
The final primary source has SHA-256
\begin{center}\footnotesize\ttfamily
79509b47cbbf20df923cc7b3e86952df33dcb3ce15b71f8e1d798212805a1902
\end{center}
The reproduction suite accepts 12 concrete closed certificates and
rejects 38 adversarial objects. Positive examples include both directions
of conditional symmetry, automatic transport at zero and nonzero
witnesses, nested iterator binders, nested paired formulas, a branch
that genuinely requires the specialized clause, and a synthesized term
containing the fresh parameter. The suite performs 12 exact translations
into the explicit-opening control, 24 exact-fuel/one-less checks, and
114 command-line checks with and without Python optimization. It also
checks a case in which outer substitution creates the naming guard:
blindly retaining the old branch fails, while the branch prescribed by
the substitution theorem succeeds.

An independent validator uses named bound variables, a separately
written sorted parser, and numerical evaluation by ordinary loops.
It passes 6,292 checks with no failures, including 700 normalization
comparisons, 1,398 numerical semantic comparisons, 500 binder-opening
comparisons, 400 simultaneous substitutions, three sets of 350
probe/clause/invariant checks, 1,500 arbitrary-JSON grammar comparisons,
and targeted scope and malformed-input examples. The source hash and
complete category counts are retained in the evidence JSON.

These executions address concrete implementation risks: capture,
accidental Python numeric coercion, malformed erased syntax, wrong
equality orientation, a nonzero reflexive base, a compound matrix,
self-containing substitutions, and eigenparameter escape. They are
finite tests. The universal soundness, substitution, completeness, and
bi-interpretation conclusions rest on the mathematical proofs and the
attributed inherited bridge, not on the number of successful tests.


\appendix
\section{Selected complete checker}
The following is the complete measured program after CPython
\texttt{ast.unparse} normalization. The archive also contains the source,
the fuel-bounded variant, its matched control, and the reproduction
script. The normalizer, all trusted syntax and scope checks, proof checking,
and the JSON interface are included in the measurement.
\lstinputlisting{code/checker.normalized.py}

\section{Reproduction and source attribution}
Run \texttt{python3 reproduce.py} from the package directory. It recomputes
source metrics and runs the delivered verification. Run
\texttt{python3 build\_pdf.py} to rebuild this PDF using an installed TeX
distribution; an ordinary \texttt{latexmk} command is also documented in
\texttt{README.txt}. Reproducing the checker validation does not require a
network connection or a proof assistant.

The \texttt{reference} directory contains attributed upstream mathematical
fragments for the native choice calculus, ordinary interpretation,
predecessor, equality reflection, and primitive-recursion compilation.
It also contains the two terminal Lean bridge modules and baseline checker
snapshots used for comparisons. These are reference inputs, not new
theorems of this package. The Lean modules import the upstream library;
the two snapshots alone are not a standalone Lean project.
\texttt{PROVENANCE.json} records exact source locations and hashes.
\texttt{CLAIMS.json} separates new theorems, inherited results, finite
implementation evidence, and claims deliberately not made.

\begin{thebibliography}{9}
\raggedright
\bibitem{repo}
ProveIt contributors, \emph{Concise presentations of intuitionistic
primitive recursive arithmetic}, README at commit
\texttt{17898ea7e6a1}.
\href{https://github.com/VladimirReshetnikov/ProveIt/tree/17898ea7e6a1300dfdcad71040d5daeed67916d4/Logic/PrimitiveRecursiveArithmetic/ConciseIPRA}{Pinned ConciseIPRA source tree}.

\bibitem{witness}
ProveIt research package, \emph{iPRA with equality at witness discharge},
\path{docs/incoming/ipra_witness_discharge_2026-11-11.zip},
at the same commit. Inherited full-normalization clause interface,
two-row equality package, and zero-witness case repair.

\bibitem{normalization}
ProveIt research package, \texttt{Research/full\_normalization/report.tex},
at the same commit. Termination, confluence, and substitution coherence
for scalar iterator normalization.

\bibitem{uniform}
ProveIt research packages, \texttt{Research/uniform\_clauses},
\texttt{Research/raw\_clausal}, and the incoming
\texttt{ipra\_uniform\_binding\_2026-11-11.zip}, at the same commit.
Clause interfaces and complete checker comparisons.

\bibitem{lean}
ProveIt contributors, \texttt{Lean/ConciseIPRA/BiInterpretation.lean}
and \texttt{NativeBiInterpretation.lean}, at the same commit.
Theorems \texttt{biInterp}, \texttt{native\_iff\_pra},
\texttt{pra\_iff\_native}, \texttt{native\_roundTrip},
\texttt{pra\_native\_roundTrip}, and
\texttt{native\_inconsistent\_iff}.

\bibitem{arithmetic}
ProveIt research package, \texttt{Research/certified\_zero/reference},
at the same commit. Native logical macros, scalar predecessor,
equality reflection, quantifier-free induction, and general
primitive-recursion compilation; reproduced reference fragments
are identified in the package provenance.
\end{thebibliography}
\end{document}
